\subsection{Lineare Ausgleichsrechnung}
Der Ansatz der Methode der kleinsten Quadrate ist in Matrixform $\displaystyle \min_{\hat{x} \in \R^n} ||A\hat{x} - b||_2^2$ und als Summe:
\drmvspace
\begin{align*}
    (a, c) = \argmin{p \in \R^n, q \in \R} \sum_{i = 1}^{m} |y_i - p^{\top} x_i - q|^2
\end{align*}

\drmvspace
Wobei $y_i$ die $y$-Koordinaten der Messpunkte zugehörig zu $x_i$ sind.

\innumpy haben wir die Funktionen \texttt{numpy.polyfit} (um ein Polynom zu fitten), oder die allgemeinere Methode \texttt{numpy.linalg.lstsq}.
Für eine Lösung mit minimaler Norm können wir die Moore-Penrose-Pseudoinverse verwenden, die \texttt{numpy.linalg.pinv} berechnet.
